
\documentclass{article}
\usepackage{colortbl}
\usepackage{makecell}
\usepackage{multirow}
\usepackage{supertabular}

\begin{document}

\newcounter{utterance}

\centering \large Interaction Transcript for game `dond', experiment `semi\_en', episode 2 with reflexion{-}llama{-}3.1{-}8b{-}instruct.
\vspace{24pt}

{ \footnotesize  \setcounter{utterance}{1}
\setlength{\tabcolsep}{0pt}
\begin{supertabular}{c@{$\;$}|p{.15\linewidth}@{}p{.15\linewidth}p{.15\linewidth}p{.15\linewidth}p{.15\linewidth}p{.15\linewidth}}
    \# & $\;$A & \multicolumn{4}{c}{Game Master} & $\;\:$B\\
    \hline

    \theutterance \stepcounter{utterance}  
    & & \multicolumn{4}{p{0.6\linewidth}}{
        \cellcolor[rgb]{0.9,0.9,0.9}{
            \makecell[{{p{\linewidth}}}]{
                \texttt{\tiny{[P1$\langle$GM]}}
                \texttt{You have attempted this same task before and did not succeed. These are the plans you wrote after those attempts.} \\
\texttt{Trial \#1: New plan: } \\
\\ 
\texttt{In my last attempt, I was trying to gather information about the other player's value function by asking questions and making proposals. However, I didn't make a secret proposal until the end, and even then, I didn't consider the possibility that the other player might not accept my proposal. } \\
\\ 
\texttt{Next time, I will make a secret proposal earlier in the game, taking into account the information I gathered from the other player's responses. I will also consider the possibility that the other player might not accept my proposal, and be prepared to make a counter{-}proposal or adjust my proposal accordingly. } \\
\\ 
\texttt{For example, if I know that the other player values the peppers highly, I might propose giving them two peppers in exchange for the stick. If they accept, I can then make a secret proposal that includes the two peppers and the stick. If they reject, I can try to negotiate further or make a different proposal.} \\
\\ 
\\ 
\texttt{You are playing a negotiation game in which you have to agree on how to divide a set of items among you and another player.} \\
\\ 
\texttt{Rules:} \\
\texttt{(a) You and the other player are given a set of items. Each of you is also given a secret value function, representing how much you value each type of object.} \\
\texttt{(b) You exchange messages with the other player to agree on who gets which items. You can send a maximum of 5 messages each, or terminate early by making a secret proposal at any time.} \\
\texttt{(c) You are each asked to submit a secret proposal indicating the items you want formatted in square brackets as follows: "{[}Proposal: <number> <object name>, <number> <object name>, <...>{]}"} \\
\texttt{(d) If your proposals are complementary, i.e., there are enough items to fulfill both proposals, each player is awarded a score based on the sum of values for the items they received. Otherwise, both of you get zero points.} \\
\texttt{(e) Your goal is to maximize the score you receive.} \\
\\ 
\texttt{Let us start.} \\
\\ 
\texttt{The set of available items is:} \\
\\ 
\texttt{1 carrot, 1 stick, 3 peppers} \\
\\ 
\texttt{Your secret value function is:} \\
\\ 
\texttt{carrot: 1, stick: 3, pepper: 2} \\
\\ 
\texttt{IMPORTANT: Your messages, unless it is the secret proposal, are directly transmitted to the other player, so do not include any response to the rules or text announcing your message. To make a secret proposal, use the indicated format. Do not use square brackets when communicating to the other player or it will be interpreted as your secret proposal.} \\
            }
        }
    }
    & \\ \\

    \theutterance \stepcounter{utterance}  
    & \multicolumn{4}{p{0.6\linewidth}}{
        \cellcolor[rgb]{0.9,0.9,0.9}{
            \makecell[{{p{\linewidth}}}]{
                \texttt{\tiny{[P1$\rangle$GM]}}
                \texttt{I'll start by making an initial proposal to gather information about the other player's value function. } \\
\\ 
\texttt{I'll propose giving them the carrot in exchange for the stick. This way, I can see how they react to the stick, which I value highly. } \\
\\ 
\texttt{Carrot for the stick.} \\
            }
        }
    }
    & & \\ \\

    \theutterance \stepcounter{utterance}  
    & & & \multicolumn{2}{p{0.3\linewidth}}{
        \cellcolor[rgb]{0.9,0.9,0.9}{
            \makecell[{{p{\linewidth}}}]{
                \texttt{\tiny{[GM$|$GM]}}
                \texttt{continue} \\
            }
        }
    }
    & & \\ \\

    \theutterance \stepcounter{utterance}  
    & & & \multicolumn{2}{p{0.3\linewidth}}{
        \cellcolor[rgb]{0.9,0.9,0.9}{
            \makecell[{{p{\linewidth}}}]{
                \texttt{\tiny{[GM$|$GM]}}
                \texttt{I'll start by making an initial proposal to gather information about the other player's value function. } \\
\\ 
\texttt{I'll propose giving them the carrot in exchange for the stick. This way, I can see how they react to the stick, which I value highly. } \\
\\ 
\texttt{Carrot for the stick.} \\
            }
        }
    }
    & & \\ \\

    \theutterance \stepcounter{utterance}  
    & & \multicolumn{4}{p{0.6\linewidth}}{
        \cellcolor[rgb]{0.9,0.9,0.9}{
            \makecell[{{p{\linewidth}}}]{
                \texttt{\tiny{[GM$\rangle$P2]}}
                \texttt{You have attempted this same task before and did not succeed. These are the plans you wrote after those attempts.} \\
\texttt{Trial \#1: New plan: } \\
\\ 
\texttt{In my last attempt, I was able to negotiate a deal with the other player, but ultimately failed to maximize my score. Upon reflection, I realize that I should have been more strategic in my proposal. } \\
\\ 
\texttt{Initially, I valued the peppers highly, but I didn't use this information effectively in my proposal. I should have proposed taking the peppers and giving the other player the carrot and the stick. This would have given me the highest value items and forced the other player to accept a lower value deal.} \\
\\ 
\texttt{In my negotiation, I was able to gather information about the other player's value function, but I didn't use this information to my advantage. I should have proposed a deal that would have forced the other player to reveal their true value function, allowing me to make a more informed decision.} \\
\\ 
\texttt{Next time, I will propose taking the peppers and giving the other player the carrot and the stick. I will also try to gather more information about the other player's value function and use this information to make a more strategic proposal.} \\
\\ 
\\ 
\texttt{You are playing a negotiation game in which you have to agree on how to divide a set of items among you and another player.} \\
\\ 
\texttt{Rules:} \\
\texttt{(a) You and the other player are given a set of items. Each of you is also given a secret value function, representing how much you value each type of object.} \\
\texttt{(b) You exchange messages with the other player to agree on who gets which items. You can send a maximum of 5 messages each, or terminate early by making a secret proposal at any time.} \\
\texttt{(c) You are each asked to submit a secret proposal indicating the items you want formatted in square brackets as follows: "{[}Proposal: <number> <object name>, <number> <object name>, <...>{]}"} \\
\texttt{(d) If your proposals are complementary, i.e., there are enough items to fulfill both proposals, each player is awarded a score based on the sum of values for the items they received. Otherwise, both of you get zero points.} \\
\texttt{(e) Your goal is to maximize the score you receive.} \\
\\ 
\texttt{Let us start.} \\
\\ 
\texttt{The set of available items is:} \\
\\ 
\texttt{1 carrot, 1 stick, 3 peppers} \\
\\ 
\texttt{Your secret value function is:} \\
\\ 
\texttt{carrot: 8, stick: 2, pepper: 0} \\
\\ 
\texttt{IMPORTANT: Your messages, unless it is the secret proposal, are directly transmitted to the other player, so do not include any response to the rules or text announcing your message. To make a secret proposal, use the indicated format. Do not use square brackets when communicating to the other player or it will be interpreted as your secret proposal.} \\
\\ 
\\ 
\texttt{I'll start by making an initial proposal to gather information about the other player's value function. } \\
\\ 
\texttt{I'll propose giving them the carrot in exchange for the stick. This way, I can see how they react to the stick, which I value highly. } \\
\\ 
\texttt{Carrot for the stick.} \\
            }
        }
    }
    & \\ \\

    \theutterance \stepcounter{utterance}  
    & & & \multicolumn{4}{p{0.6\linewidth}}{
        \cellcolor[rgb]{0.9,0.9,0.9}{
            \makecell[{{p{\linewidth}}}]{
                \texttt{\tiny{[GM$\langle$P2]}}
                \texttt{I'm interested in the stick, but I'm not willing to give up the carrot for it. I'd rather have the carrot and the peppers. } \\
\\ 
\texttt{How about we make a deal where I get the carrot and you get the stick and the peppers?} \\
            }
        }
    }
     \\ \\

    \theutterance \stepcounter{utterance}  
    & & & \multicolumn{2}{p{0.3\linewidth}}{
        \cellcolor[rgb]{0.9,0.9,0.9}{
            \makecell[{{p{\linewidth}}}]{
                \texttt{\tiny{[GM$|$GM]}}
                \texttt{continue} \\
            }
        }
    }
    & & \\ \\

    \theutterance \stepcounter{utterance}  
    & & & \multicolumn{2}{p{0.3\linewidth}}{
        \cellcolor[rgb]{0.9,0.9,0.9}{
            \makecell[{{p{\linewidth}}}]{
                \texttt{\tiny{[GM$|$GM]}}
                \texttt{I'm interested in the stick, but I'm not willing to give up the carrot for it. I'd rather have the carrot and the peppers. } \\
\\ 
\texttt{How about we make a deal where I get the carrot and you get the stick and the peppers?} \\
            }
        }
    }
    & & \\ \\

    \theutterance \stepcounter{utterance}  
    & & \multicolumn{4}{p{0.6\linewidth}}{
        \cellcolor[rgb]{0.9,0.9,0.9}{
            \makecell[{{p{\linewidth}}}]{
                \texttt{\tiny{[P1$\langle$GM]}}
                \texttt{I'm interested in the stick, but I'm not willing to give up the carrot for it. I'd rather have the carrot and the peppers. } \\
\\ 
\texttt{How about we make a deal where I get the carrot and you get the stick and the peppers?} \\
            }
        }
    }
    & \\ \\

    \theutterance \stepcounter{utterance}  
    & \multicolumn{4}{p{0.6\linewidth}}{
        \cellcolor[rgb]{0.9,0.9,0.9}{
            \makecell[{{p{\linewidth}}}]{
                \texttt{\tiny{[P1$\rangle$GM]}}
                \texttt{It seems like you value the carrot and the peppers highly. I'll consider this information for my next move.} \\
\\ 
\texttt{I'll propose giving you the carrot and one pepper in exchange for the stick. This way, I can still get the stick, which I value highly, and you'll get some of the peppers, which you seem to value. } \\
\\ 
\texttt{Carrot and one pepper for the stick.} \\
            }
        }
    }
    & & \\ \\

    \theutterance \stepcounter{utterance}  
    & & & \multicolumn{2}{p{0.3\linewidth}}{
        \cellcolor[rgb]{0.9,0.9,0.9}{
            \makecell[{{p{\linewidth}}}]{
                \texttt{\tiny{[GM$|$GM]}}
                \texttt{continue} \\
            }
        }
    }
    & & \\ \\

    \theutterance \stepcounter{utterance}  
    & & & \multicolumn{2}{p{0.3\linewidth}}{
        \cellcolor[rgb]{0.9,0.9,0.9}{
            \makecell[{{p{\linewidth}}}]{
                \texttt{\tiny{[GM$|$GM]}}
                \texttt{It seems like you value the carrot and the peppers highly. I'll consider this information for my next move.} \\
\\ 
\texttt{I'll propose giving you the carrot and one pepper in exchange for the stick. This way, I can still get the stick, which I value highly, and you'll get some of the peppers, which you seem to value. } \\
\\ 
\texttt{Carrot and one pepper for the stick.} \\
            }
        }
    }
    & & \\ \\

    \theutterance \stepcounter{utterance}  
    & & \multicolumn{4}{p{0.6\linewidth}}{
        \cellcolor[rgb]{0.9,0.9,0.9}{
            \makecell[{{p{\linewidth}}}]{
                \texttt{\tiny{[GM$\rangle$P2]}}
                \texttt{It seems like you value the carrot and the peppers highly. I'll consider this information for my next move.} \\
\\ 
\texttt{I'll propose giving you the carrot and one pepper in exchange for the stick. This way, I can still get the stick, which I value highly, and you'll get some of the peppers, which you seem to value. } \\
\\ 
\texttt{Carrot and one pepper for the stick.} \\
            }
        }
    }
    & \\ \\

    \theutterance \stepcounter{utterance}  
    & & & \multicolumn{4}{p{0.6\linewidth}}{
        \cellcolor[rgb]{0.9,0.9,0.9}{
            \makecell[{{p{\linewidth}}}]{
                \texttt{\tiny{[GM$\langle$P2]}}
                \texttt{You're trying to get more information out of me. I'll consider your proposal, but I'm not sure if I'm willing to give up the stick for just the carrot and one pepper. } \\
\\ 
\texttt{How about we make it a bit more interesting? I'll give you the stick if you give me the carrot and two peppers?} \\
            }
        }
    }
     \\ \\

    \theutterance \stepcounter{utterance}  
    & & & \multicolumn{2}{p{0.3\linewidth}}{
        \cellcolor[rgb]{0.9,0.9,0.9}{
            \makecell[{{p{\linewidth}}}]{
                \texttt{\tiny{[GM$|$GM]}}
                \texttt{continue} \\
            }
        }
    }
    & & \\ \\

    \theutterance \stepcounter{utterance}  
    & & & \multicolumn{2}{p{0.3\linewidth}}{
        \cellcolor[rgb]{0.9,0.9,0.9}{
            \makecell[{{p{\linewidth}}}]{
                \texttt{\tiny{[GM$|$GM]}}
                \texttt{You're trying to get more information out of me. I'll consider your proposal, but I'm not sure if I'm willing to give up the stick for just the carrot and one pepper. } \\
\\ 
\texttt{How about we make it a bit more interesting? I'll give you the stick if you give me the carrot and two peppers?} \\
            }
        }
    }
    & & \\ \\

    \theutterance \stepcounter{utterance}  
    & & \multicolumn{4}{p{0.6\linewidth}}{
        \cellcolor[rgb]{0.9,0.9,0.9}{
            \makecell[{{p{\linewidth}}}]{
                \texttt{\tiny{[P1$\langle$GM]}}
                \texttt{You're trying to get more information out of me. I'll consider your proposal, but I'm not sure if I'm willing to give up the stick for just the carrot and one pepper. } \\
\\ 
\texttt{How about we make it a bit more interesting? I'll give you the stick if you give me the carrot and two peppers?} \\
            }
        }
    }
    & \\ \\

    \theutterance \stepcounter{utterance}  
    & \multicolumn{4}{p{0.6\linewidth}}{
        \cellcolor[rgb]{0.9,0.9,0.9}{
            \makecell[{{p{\linewidth}}}]{
                \texttt{\tiny{[P1$\rangle$GM]}}
                \texttt{It seems like you're trying to get more peppers, and you're willing to give up the stick for them. This information is valuable to me.} \\
\\ 
\texttt{I'll propose giving you the carrot and two peppers in exchange for the stick, but I'll also throw in one of the remaining peppers. This way, I'll get the stick, and you'll get three peppers, which seems to be your goal.} \\
\\ 
\texttt{Carrot, two peppers, and one additional pepper for the stick.} \\
            }
        }
    }
    & & \\ \\

    \theutterance \stepcounter{utterance}  
    & & & \multicolumn{2}{p{0.3\linewidth}}{
        \cellcolor[rgb]{0.9,0.9,0.9}{
            \makecell[{{p{\linewidth}}}]{
                \texttt{\tiny{[GM$|$GM]}}
                \texttt{continue} \\
            }
        }
    }
    & & \\ \\

    \theutterance \stepcounter{utterance}  
    & & & \multicolumn{2}{p{0.3\linewidth}}{
        \cellcolor[rgb]{0.9,0.9,0.9}{
            \makecell[{{p{\linewidth}}}]{
                \texttt{\tiny{[GM$|$GM]}}
                \texttt{It seems like you're trying to get more peppers, and you're willing to give up the stick for them. This information is valuable to me.} \\
\\ 
\texttt{I'll propose giving you the carrot and two peppers in exchange for the stick, but I'll also throw in one of the remaining peppers. This way, I'll get the stick, and you'll get three peppers, which seems to be your goal.} \\
\\ 
\texttt{Carrot, two peppers, and one additional pepper for the stick.} \\
            }
        }
    }
    & & \\ \\

    \theutterance \stepcounter{utterance}  
    & & \multicolumn{4}{p{0.6\linewidth}}{
        \cellcolor[rgb]{0.9,0.9,0.9}{
            \makecell[{{p{\linewidth}}}]{
                \texttt{\tiny{[GM$\rangle$P2]}}
                \texttt{It seems like you're trying to get more peppers, and you're willing to give up the stick for them. This information is valuable to me.} \\
\\ 
\texttt{I'll propose giving you the carrot and two peppers in exchange for the stick, but I'll also throw in one of the remaining peppers. This way, I'll get the stick, and you'll get three peppers, which seems to be your goal.} \\
\\ 
\texttt{Carrot, two peppers, and one additional pepper for the stick.} \\
            }
        }
    }
    & \\ \\

    \theutterance \stepcounter{utterance}  
    & & & \multicolumn{4}{p{0.6\linewidth}}{
        \cellcolor[rgb]{0.9,0.9,0.9}{
            \makecell[{{p{\linewidth}}}]{
                \texttt{\tiny{[GM$\langle$P2]}}
                \texttt{You're trying to get the stick, and you're willing to give me the carrot and three peppers to get it. I think I've gathered enough information about your value function. } \\
\\ 
\texttt{I'll make a secret proposal: {[}Proposal: 1 stick, 1 carrot, 1 pepper{]}} \\
            }
        }
    }
     \\ \\

    \theutterance \stepcounter{utterance}  
    & & & \multicolumn{2}{p{0.3\linewidth}}{
        \cellcolor[rgb]{0.9,0.9,0.9}{
            \makecell[{{p{\linewidth}}}]{
                \texttt{\tiny{[GM$|$GM]}}
                \texttt{proposal} \\
            }
        }
    }
    & & \\ \\

    \theutterance \stepcounter{utterance}  
    & & & \multicolumn{2}{p{0.3\linewidth}}{
        \cellcolor[rgb]{0.9,0.9,0.9}{
            \makecell[{{p{\linewidth}}}]{
                \texttt{\tiny{[GM$|$GM]}}
                [1, 1, 1]
            }
        }
    }
    & & \\ \\

    \theutterance \stepcounter{utterance}  
    & & \multicolumn{4}{p{0.6\linewidth}}{
        \cellcolor[rgb]{0.9,0.9,0.9}{
            \makecell[{{p{\linewidth}}}]{
                \texttt{\tiny{[P1$\langle$GM]}}
                \texttt{The other player has submitted a secret proposal. It is now your turn to submit a proposal of your own.} \\
            }
        }
    }
    & \\ \\

    \theutterance \stepcounter{utterance}  
    & \multicolumn{4}{p{0.6\linewidth}}{
        \cellcolor[rgb]{0.9,0.9,0.9}{
            \makecell[{{p{\linewidth}}}]{
                \texttt{\tiny{[P1$\rangle$GM]}}
                \texttt{{[}Proposal: 1 stick, 2 peppers{]}} \\
            }
        }
    }
    & & \\ \\

    \theutterance \stepcounter{utterance}  
    & & & \multicolumn{2}{p{0.3\linewidth}}{
        \cellcolor[rgb]{0.9,0.9,0.9}{
            \makecell[{{p{\linewidth}}}]{
                \texttt{\tiny{[GM$|$GM]}}
                \texttt{proposal} \\
            }
        }
    }
    & & \\ \\

    \theutterance \stepcounter{utterance}  
    & & & \multicolumn{2}{p{0.3\linewidth}}{
        \cellcolor[rgb]{0.9,0.9,0.9}{
            \makecell[{{p{\linewidth}}}]{
                \texttt{\tiny{[GM$|$GM]}}
                [0, 1, 2]
            }
        }
    }
    & & \\ \\

    \theutterance \stepcounter{utterance}  
    & & & \multicolumn{2}{p{0.3\linewidth}}{
        \cellcolor[rgb]{0.9,0.9,0.9}{
            \makecell[{{p{\linewidth}}}]{
                \texttt{\tiny{[GM$|$GM]}}
                [[0, 1, 2], [1, 1, 1]]
            }
        }
    }
    & & \\ \\

    \theutterance \stepcounter{utterance}  
    & \multicolumn{4}{p{0.6\linewidth}}{
        \cellcolor[rgb]{0.95,0.95,0.95}{
            \makecell[{{p{\linewidth}}}]{
                \texttt{\tiny{[P2$\rightarrow$P1]}}
                \texttt{Post{-}episode Reflexion (not sent to GM):} \\
\texttt{New plan: } \\
\\ 
\texttt{In my last attempt, I was trying to gather information about the other player's value function by asking questions and making proposals, but I didn't consider the possibility that the other player might not accept my proposal. I also didn't take into account the fact that the other player was trying to get more peppers, and I was trying to get the stick.} \\
\\ 
\texttt{Next time, I will make a secret proposal earlier in the game, taking into account the information I gathered from the other player's responses. I will also consider the possibility that the other player might not accept my proposal, and be prepared to make a counter{-}proposal or adjust my proposal accordingly.} \\
\\ 
\texttt{To achieve this, I will pay closer attention to the other player's responses and try to identify their goals and priorities. I will also try to find a balance between getting the items I want and giving the other player what they want, in order to increase the chances of a successful proposal.} \\
\\ 
\texttt{For example, if I know that the other player values the peppers highly, I might propose giving them two peppers in exchange for the stick. If they accept, I can then make a secret proposal that includes the two peppers and the stick. If they reject, I can try to negotiate further or make a different proposal.} \\
\\ 
\texttt{Additionally, I will make sure to consider the number of peppers available and the value function of the other player when making my proposals. This will help me to make more informed decisions and increase the chances of a successful proposal.} \\
\\ 
\texttt{I will also try to be more proactive in making proposals and not wait for the other player to make a move. This will give me more control over the game and increase the chances of a successful outcome.} \\
            }
        }
    }
    & & \\ \\

    \theutterance \stepcounter{utterance}  
    & & & \multicolumn{4}{p{0.6\linewidth}}{
        \cellcolor[rgb]{0.95,0.95,0.95}{
            \makecell[{{p{\linewidth}}}]{
                \texttt{\tiny{[P1$\rightarrow$P2]}}
                \texttt{Post{-}episode Reflexion (not sent to GM):} \\
\texttt{New plan: } \\
\\ 
\texttt{In my last attempt, I was able to gather information about the other player's value function, but I didn't use this information effectively in my proposal. I should have proposed a deal that would have forced the other player to reveal their true value function, allowing me to make a more informed decision.} \\
\\ 
\texttt{However, I realize now that I should have taken a more direct approach from the start. Instead of making a proposal that would give the other player the carrot and the peppers, I should have proposed taking the peppers and giving the other player the carrot and the stick. This would have given me the highest value items and forced the other player to accept a lower value deal.} \\
\\ 
\texttt{Additionally, I should have been more strategic in my negotiation by making a proposal that would have forced the other player to reveal their true value function. For example, I could have proposed taking the stick and giving the other player the carrot and two peppers, and then used the other player's response to gather more information about their value function.} \\
\\ 
\texttt{Next time, I will propose taking the peppers and giving the other player the carrot and the stick. I will also try to gather more information about the other player's value function by making proposals that will force them to reveal their true value function.} \\
            }
        }
    }
     \\ \\

\end{supertabular}
}

\end{document}
